\SourceLeaf{207}
dort rechts liegt, was hier links ist. Dieser Unterschied ist aber kein wesentlicher, und verschwindet, wenn man in der einen Ebne diejenige Seite, welche man vorher als obere betrachtete, zur untern macht. Diese letztre Bemerkung lässt sich übrigens allemal in Anwendung bringen, wenn die eine der beiden Flächen eine Ebne ist, daher wir in den folgenden Beispielen dieser Art uns bloss auf die erste Auflösung beschränken können.

\Article{9}
Wir wollen nun (als zweites Beispiel) die Darstellung der Fläche eines geraden Kegels in der Ebne betrachten. Als Gleichung der erstern nehmen wir an
\[
xx+yy-kkzz=0
\]
\noindent wo wir ferner
\[
\begin{aligned}
x&=kt\cos u\\
y&=kt\sin u\\
z&=t
\end{aligned}
\]
\noindent und wie vorhin \(X=T,Y=U,Z=0\) setzen.

Die Differentialgleichung
\[
\omega=(kk+1)\mathrm{d}t^2+kktt\,\mathrm{d}u^2=0
\]
\noindent gibt hier die beiden Integrale
\[
\log t\pm i\BareRoot{\frac{kk}{kk+1}}\mathord{.}u=\text{Const.}
\]
\noindent Wir haben demnach die Auflösung
\[
\begin{aligned}
X+iY&=f(\log t+i\BareRoot{\frac{kk}{kk+1}}\mathord{.}u)\\
X-iY&=f'(\log t-i\BareRoot{\frac{kk}{kk+1}}\mathord{.}u)
\end{aligned}
\]
\noindent d. i. es wird, indem \(f\) eine willkürliche Function bedeutet, für \(X\) der reelle Theil von
\[
f(\log t+i\BareRoot{\frac{kk}{kk+1}}\mathord{.}u)
\]
\noindent und für \(Y\) der imaginäre, nach Weglassung des Factors \(i\), angenommen.
